% ==============================================================================
% Formation C# & SQL — Module 001 : Les Bases (Module Complémentaire / Bonus)
% Fichier : src/P1_02_b.tex
% Sujet   : Conventions de Nommage, Formatage, Clean Code et Outillage C#
% Note    : Pas de page de titre (suite directe de P1_01_a.tex)
% ==============================================================================

\section{Conventions de Nommage, Formatage et Clean Code}
\marginpar{\tiny\color{gray}P1\_02\_b.tex}

Écrire du code informatique ne se limite pas à produire une séquence d'instructions exécutables par le compilateur. Dans l'ingénierie logicielle professionnelle, le code C\# doit être \textbf{lisible}, \textbf{maintenable} et \textbf{homogène}. Les conventions de style constituent le langage commun facilitant la collaboration au sein des équipes de développement.



\subsection{Conventions de Nommage Officielles}

Le nommage d'un élément logiciel doit rendre le code auto-documenté. Microsoft définit des règles précises de casse selon la nature et la portée de chaque symbole.

\subsubsection{Les Styles de Casse en C\#}

\begin{center}
    \textbf{Tableau : Règles de casse officielles C\# (Microsoft Design Guidelines)} \\[0.3cm]
    \begin{tabular}{@{}lllp{6cm}@{}}
        \toprule
        \textbf{Style} & \textbf{Format} & \textbf{Exemple} & \textbf{Cas d'utilisation} \\
        \midrule
        \textbf{PascalCase} & Majuscule à chaque mot & \texttt{CalculerPrixTotal} & Classes, structures, enums, méthodes, propriétés, événements, constantes (\texttt{const} / \texttt{readonly}). \\
        \textbf{camelCase} & 1\textsuperscript{er} mot en minuscule & \texttt{prixUnitaire} & Variables locales, paramètres de méthodes. \\
        \textbf{\_camelCase} & camelCase précédé de \texttt{\_} & \texttt{\_compteurClient} & Champs privés d'une classe (\emph{private fields}). \\
        \textbf{IPascalCase} & PascalCase avec préfixe \texttt{I} & \texttt{IDisposable} & Interfaces. \\
        \bottomrule
    \end{tabular}
\end{center}

\subsubsection{Exemples de Conformité Syntaxique}

\begin{lstlisting}[caption={Conventions de nommage recommandées en C\#}]
// Interface : Prefix 'I' + PascalCase
public interface ICalculable
{
    // Propriete : PascalCase
    decimal Total { get; }
}

// Classe : PascalCase
public class GestionnaireCommande : ICalculable
{
    // Constante : PascalCase
    public const int LimiteArticles = 100;

    // Champ prive : _camelCase
    private readonly string _identifiantCommande;
    private int _nombreArticles;

    // Propriete publique : PascalCase
    public decimal Total { get; private set; }

    // Constructeur : parametre en camelCase
    public GestionnaireCommande(string identifiantCommande)
    {
        _identifiantCommande = identifiantCommande;
    }

    // Methode : PascalCase, parametre : camelCase, variable locale : camelCase
    public void AjouterArticle(decimal prixUnitaire, int quantite)
    {
        decimal sousTotal = prixUnitaire * quantite;
        Total += sousTotal;
        _nombreArticles += quantite;
    }
}
\end{lstlisting}

\begin{tcolorbox}[colback=backcolour,colframe=codegray,title=\textbf{Pièges de nommage à éviter}]
    \begin{itemize}
        \item \textbf{Ne pas utiliser le \texttt{snake\_case}} (\texttt{prix\_total}) : réservé à d'autres langages (Python, Rust).
        \item \textbf{Éviter la notation hongroise} (\texttt{strNom}, \texttt{intAge}) : le typage C\# est déjà statique et explicite.
        \item \textbf{Omettre le préfixe \texttt{I} sur les interfaces} : rend la distinction impossible avec les classes abstraites.
    \end{itemize}
\end{tcolorbox}



\subsection{Conventions de Formatage (Style Allman)}

Le formatage régit l'organisation visuelle du code source. C\# adopte historiquement le style d'accolades \textbf{Allman}.

\subsubsection{Règles Fondamentales d'Organigramme Visuel}

\begin{enumerate}
    \item \textbf{Accolades Allman} : L'accolade ouvrante \texttt{\{} se place \textbf{toujours sur sa propre ligne}, alignée verticalement avec la structure parente.
    \item \textbf{Indentation} : Strictement 4 espaces par niveau de bloc (ne pas mélanger tabulations et espaces).
    \item \textbf{Espacement des opérateurs} : Un espace unique de part et d'autre des opérateurs binaires (\texttt{=}, \texttt{+}, \texttt{==}, \texttt{\&\&}).
\end{enumerate}

\begin{lstlisting}[caption={Comparaison de formatage : Allman (C\#) vs K\&R (Java/JS)}]
// [CORRECT] Style Allman officiel C#
public void InscrireUtilisateur(string nom)
{
    if (string.IsNullOrWhiteSpace(nom))
    {
        throw new ArgumentException("Nom invalide", nameof(nom));
    }
}

// [INCORRECT] Style K&R (Anti-pattern en C#)
public void InscrireUtilisateur(string nom) {
    if (string.IsNullOrWhiteSpace(nom)) {
        throw new ArgumentException("Nom invalide", nameof(nom));
    }
}
\end{lstlisting}



\subsection{Organisation Interne d'une Classe C\#}

Une classe lisible regroupe ses membres selon un ordre standardisé reconnu par la communauté.

\begin{center}
    \textbf{Tableau : Ordre d'agencement recommandé des membres d'une classe} \\[0.3cm]
    \begin{tabular}{@{}c>{\bfseries}lp{8.5cm}@{}}
        \toprule
        \textbf{Ordre} & \textbf{Catégorie de Membre} & \textbf{Description} \\
        \midrule
        1 & Constantes & \texttt{public const} et \texttt{private const}. \\
        2 & Champs privés & Champs d'instance d'encapsulation (\texttt{\_camelCase}). \\
        3 & Constructeurs & Initialisateurs d'instance. \\
        4 & Propriétés publiques & Accesseurs de données (\texttt{get}, \texttt{set}, \texttt{init}). \\
        5 & Méthodes publiques & Points d'entrée de la logique métier. \\
        6 & Méthodes privées & Méthodes utilitaires internes d'assistance. \\
        \bottomrule
    \end{tabular}
\end{center}

\begin{lstlisting}[caption={Exemple de structure canonique d'un service C\#}]
namespace LesBases.Services;

public class ServiceCompteBancaire
{
    // 1. Constantes
    private const decimal SoldeMinimumAutorise = -500.00m;

    // 2. Champs privés
    private readonly string _numeroCompte;
    private decimal _solde;

    // 3. Constructeurs
    public ServiceCompteBancaire(string numeroCompte, decimal soldeInitial)
    {
        _numeroCompte = numeroCompte;
        _solde = soldeInitial;
    }

    // 4. Proprietes publiques
    public string NumeroCompte => _numeroCompte;
    public decimal Solde => _solde;

    // 5. Methodes publiques
    public void Deposer(decimal montant)
    {
        ValiderMontantPositif(montant);
        _solde += montant;
    }

    // 6. Methodes privees d'assistance
    private static void ValiderMontantPositif(decimal montant)
    {
        if (montant <= 0)
        {
            throw new ArgumentOutOfRangeException(nameof(montant), "Le montant doit etre positif.");
        }
    }
}
\end{lstlisting}



\subsection{Inférence de Type (\texttt{var}) vs Type Explicite}

Le mot-clé \texttt{var} (introduit en C\# 3) permet l'inférence de type à la compilation. Son utilisation est encadrée par des critères d'évidence visuelle.

\begin{tcolorbox}[colback=white,colframe=keywordcolor,title=\textbf{Règle d'or de l'inférence \texttt{var}}]
    Utiliser \texttt{var} uniquement lorsque le type retourné est \textbf{immédiatement explicite} à la lecture du côté droit de l'affectation (opérateur \texttt{new} ou littéral). Dans le cas contraire, conserver le type explicite pour maintenir la clarté.
\end{tcolorbox}

\begin{lstlisting}[caption={Usage correct et incorrect de var}]
// [BON] Le type est evident via l'instanciation directe ou le litteral
var nomUtilisateur = "Alice";                             // string
var listeClients = new List<string>();                   // List<string>
var dictionnaire = new Dictionary<int, string>();         // Evite la repetition inutile

// [MAUVAIS] Le type est masque par l'appel de methode
var resultat = CalculerIndiceForme();                     // De quel type est resultat ? int ? double ? CustomType ?

// [CORRECT] Type explicite quand l'intention doit être clarifiée
decimal indiceForme = CalculerIndiceForme();
\end{lstlisting}



\subsection{Contrôle de Flux et Retours Anticipés (\emph{Early Return})}

Pour éviter l'empilement excessif de conditions imbriquées (syndrome de la pyramide d'indentation), C\# préconise la technique des \textbf{retours anticipés}.

\begin{lstlisting}[caption={Réfractorisation d'une méthode complexe par Early Return}]
// [MAUVAIS] Imbrication profonde et illisible
public string TraiterCommande(Commande cmd)
{
    if (cmd != null)
    {
        if (cmd.EstValidee)
        {
            if (cmd.StockDisponible)
            {
                return "Commande expediee";
            }
            else
            {
                return "Rupture de stock";
            }
        }
        else
        {
            return "Commande non validee";
        }
    }
    return "Erreur commande nulle";
}

// [EXCELLENT] Structure avec retours anticipes (Early Return)
public string TraiterCommandePropre(Commande cmd)
{
    if (cmd == null) return "Erreur commande nulle";
    if (!cmd.EstValidee) return "Commande non validee";
    if (!cmd.StockDisponible) return "Rupture de stock";

    return "Commande expediee";
}
\end{lstlisting}



\subsection{Gestion Rigoureuse des Exceptions}

La gestion des erreurs doit respecter la chaîne de rétropropagation des exceptions sans altérer la pile d'appels (\emph{stack trace}).

\begin{tcolorbox}[colback=backcolour,colframe=codegray,title=\textbf{Règle d'or : \texttt{throw;} vs \texttt{throw ex;}}]
    Lors du re-lancement d'une exception dans un bloc \texttt{catch}, utiliser strictement \textbf{\texttt{throw;}} sans argument. L'instruction \texttt{throw ex;} réinitialise la pile d'appels au niveau du catch, masquant l'origine exacte du plantage.
\end{tcolorbox}

\begin{lstlisting}[caption={Bonnes pratiques de re-lancement d'exceptions}]
try
{
    ExecuterTraitementReseau();
}
catch (TimeoutException ex)
{
    // [BON] Conserve l'integralite de la stack trace originale
    Logger.LogError(ex, "Timeout reseau detecte");
    throw; 
}
catch (Exception)
{
    // [INCORRECT] Ecraser la stack trace originale
    // throw ex; <-- NE JAMAIS FAIRE CELA
}
\end{lstlisting}



\subsection{Outillage Automatisé (\texttt{.editorconfig} \& Roslyn)}

Afin de garantir le respect des conventions sans effort manuel, l'écosystème .NET utilise des fichiers de configuration standardisés à la racine du dépôt.

\subsubsection{Fichier de Configuration \texttt{.editorconfig}}

\begin{lstlisting}[caption={Exemple de fichier .editorconfig standard pour projet C\#}]
root = true

[*.cs]
# Indentation
indent_style = space
indent_size = 4

# Regles de nommage des champs prives
dotnet_naming_rule.private_fields_should_have_underscore.severity = warning
dotnet_naming_rule.private_fields_should_have_underscore.symbols  = private_fields
dotnet_naming_rule.private_fields_should_have_underscore.style    = prefix_underscore

dotnet_naming_symbols.private_fields.applicable_kinds = field
dotnet_naming_symbols.private_fields.applicable_accessibilities = private

dotnet_naming_style.prefix_underscore.required_prefix = _
dotnet_naming_style.prefix_underscore.capitalization = camel_case

# Respect obligatoire des accolades Allman
csharp_prefer_braces = true:warning
\end{lstlisting}

\subsubsection{Commande de Formatage CLI}

La CLI .NET intègre un outil natif pour corriger le formatage du projet entier :

\begin{lstlisting}[caption={Commandes dotnet format}]
# Reformate automatiquement tous les fichiers de la solution selon le .editorconfig
dotnet format

# Verification sans modification (ideal pour les pipelines CI/CD)
dotnet format --verify-no-changes
\end{lstlisting}



\subsection{Synthèse des Conventions}

\begin{center}
    \textbf{Tableau : Tableau récapitulatif des règles d'ingénierie C\#} \\[0.3cm]
    \begin{tabular}{@{}lll@{}}
        \toprule
        \textbf{Élément} & \textbf{Convention Officielle} & \textbf{Exemple Validé} \\
        \midrule
        Classe / Structure & PascalCase & \texttt{class CompteUtilisateur} \\
        Interface & I + PascalCase & \texttt{interface IRepository} \\
        Méthode / Propriété & PascalCase & \texttt{public void Sauvegarder()} \\
        Variable locale & camelCase & \texttt{int totalElements = 0;} \\
        Champ privé & \_camelCase & \texttt{private string \_cleApi;} \\
        Accolades & Style Allman & Première ligne dédiée \\
        Re-throw exception & \texttt{throw;} & Préserve la pile d'appels \\
        \bottomrule
    \end{tabular}
\end{center}
